\section{Related Work}
\label{sec:related_work}

The literature on model merging has expanded rapidly as deep neural networks have scaled in size. In this section, we situate \textbf{FluidMerge} within the landscape of flat Euclidean model merging, geometry-constrained merging, test-time adaptation, and the integration of continuous physical systems in deep learning.

\subsection{Flat Euclidean Model Merging}
Early and foundational approaches to model merging operate under the implicit assumption of a flat Euclidean parameter space. \citet{ilharco2023editing} introduced \emph{Task Arithmetic}, which constructs task-specific parameter vectors (the difference between fine-tuned and pretrained weights) and linearly adds them to the pretrained base. While simple and elegant, linear addition frequently triggers catastrophic task interference when weights conflict. To address this, \citet{yadav2024ties} developed \emph{Ties-Merging}, which resolves sign conflicts and prunes low-magnitude updates, while \citet{yu2023language} proposed \emph{DARE} to randomly drop fine-tuned weights and rescale the remaining ones to match the original activation distribution. Other methods, such as \emph{RegMean} \cite{jin2023raw}, compute closed-form linear projections based on input covariance matrices. Despite their empirical utility, all these methods rely on static, coordinate-wise operations that ignore the non-convex geometry of neural loss landscapes, leading to suboptimal performance when task representations collide \cite{yang2024adamerging, chen2024survey}.

\subsection{Manifold and Geometry-Constrained Merging}
Recognizing the limitations of flat Euclidean spaces, several researchers have turned to geometry and manifold theory to align neural networks before or during merging. \citet{yang2026orthomerge} proposed \emph{OrthoMerge}, which maps model parameters into the Lie algebra of orthogonal transformations, performs averaging under magnitude correction, and maps them back. Quotient space approaches such as \emph{GeoMerge} \cite{ainsworth2022git} leverage permutation and scaling symmetries to merge models on a quotient manifold, while \citet{peebles2020hessian} and \emph{EpiMer} recast merging as finding the Fr\'echet mean under curvature metrics linked to the expected Hessian. 

Importantly, our continuous-time weight trajectory shares deep mathematical connections with optimal transport-based model merging, specifically \emph{OT-Fusion} \cite{singh2020model}. In optimal transport, the celebrated \emph{Benamou-Brenier formulation} \cite{benamou2000formula} recharacterizes static Monge-Kantorovich transport as a continuous-time, dynamic fluid flow minimizing kinetic energy subject to a conservation-of-mass continuity equation. FluidMerge can be viewed as a parameter-space analog of such Wasserstein gradient flows: instead of transporting probability mass in activation space, we transport representation coordinates through the non-convex weight landscape, where self-supervised advection acts as the driving velocity field and Fisher-Information viscosity acts as the regularizing metric. Although prior methods acknowledge the curved geometry of deep networks, they are generally limited to static geodesic interpolations or computationally heavy algebraic operations (e.g., SVD scaling at $O(d^3)$ in OrthoMerge) rather than solving a continuous-time physical trajectory.

\subsection{Test-Time Adaptation and Self-Supervised Merging}
Our work is closely related to methods that adapt merging parameters at test time using unlabeled data. \citet{yang2024adamerging} introduced \emph{AdaMerging}, which dynamically learns merging coefficients for each layer using entropy minimization or unlabeled test data. \citet{jung2025symerge} proposed \emph{SyMerge}, which uses teacher-guided self-labeling to adapt a single task-specific layer alongside merging coefficients on unlabeled test data, mitigating the instability of entropy minimization. Additionally, \emph{Task Surgery} \cite{zhang2024task} utilizes L1-based feature alignment to isolate and merge orthogonal task representations. While these techniques introduce dynamic adjustment, they either restrict adaptation to a tiny subset of parameters (like single-layer tuning in SyMerge) or rely on static weight scaling. FluidMerge extends this line of research by simulating a continuous-time advection-diffusion trajectory over the entire parameter space of the image encoder, exposing a fundamental representation domain shift barrier that limits post-hoc self-training.

\subsection{Continuous Physical Systems in Deep Learning}
The integration of physical metaphors and continuous-time ordinary differential equations (ODEs) into deep learning was popularized by \emph{Neural ODEs} \cite{chen2018neural}, which represent the forward pass of residual networks as continuous-time activation trajectories. Similarly, Physics-Informed Neural Networks (PINNs) \cite{raissi2019physics} integrate physical conservation laws directly into neural training. Beyond activations, researchers have explored continuous trajectories for optimizer design \cite{muehlebach2021continuous} and generative modeling via diffusion processes \cite{ho2020denoising, song2020score}. However, while physical equations have been widely applied to model activations, gradients, or data distributions, \textbf{FluidMerge} is the first to model the trajectory of the \emph{network weights themselves} as a continuous, viscous fluid flow. This conceptual leap opens an entirely new category of physical, non-linear parameter blending.
